% ============================================================================
% RELATED: NEURO-VECTOR-SYMBOLIC ARCHITECTURES AND HD COMPUTING HARDWARE
% ============================================================================
\subsection{Neuro-Vector-Symbolic Architectures and HD Computing Hardware}
\label{sec:related_nvsa}

\paragraph{The NVSA Bridge.}
Hersche et al.\ \citep{hersche2023nvsa} demonstrated that combining
neural front-ends with VSA back-ends solves Raven's Progressive
Matrices---a benchmark for abstract visual reasoning---more accurately
and with orders-of-magnitude fewer gradient steps than pure deep-learning
baselines.  Their key insight is that the VSA's algebraic closure under
binding and bundling provides a native factorization substrate: the
network learns to project perception into a space where combinatorial
decomposition is a single algebraic operation.

The Value architecture pushes this factorization principle in a specific
direction: it eliminates the neural front-end entirely.  Where NVSA uses
a learned encoder to project raw input into VSA space, the Value
encoding assigns deterministic binary representations directly from raw
bytes.  This removes encoder training at the cost of a rigid,
input-dependent encoding that cannot adapt to the statistics of a
particular dataset.  The trade-off is clear: guaranteed determinism and
reproducibility versus the flexibility of learned representations.

\paragraph{HD Computing as a Hardware Paradigm.}
Rahimi et al.\ \citep{rahimi2017hd} established hyperdimensional
computing as a nanoscalable paradigm, demonstrating that HD vectors with
thousands of dimensions can perform classification, language recognition,
and biosignal processing with fast-learning algorithms amenable to 3D
nanometer circuit technology.  Their architecture exploits concentration
of measure in high-dimensional spaces: randomly sampled HD vectors are
approximately orthogonal, enabling robust distributed representations.

The Value operates at a dimensionality ($D = 8192$) within the classical
HD regime ($D \ge 4{,}000$), gaining the concentration-of-measure
benefits of high-dimensional spaces.  The compensation for any residual
collision risk between similar inputs is the population of frames in
the runtime: rather than encoding all information into a single
high-dimensional vector, the system distributes content across many
frames that interact through Boolean operations and, after pruning,
through PGA transition motors.

\paragraph{VSA on Emerging Hardware.}
Kleyko et al.\ \citep{kleyko2022vsa} survey VSA as a computing framework
for emerging hardware, from in-memory computing to neuromorphic chips.
They identify the core VSA operations---binding, bundling, and
permutation---as naturally parallelizable primitives that map efficiently
to non-von-Neumann architectures.  The multi-substrate backend described
in \Cref{sec:multi_backend} instantiates this principle: the entire
filtering pipeline reduces to SIMD bitwise operations with no
data-dependent branching, making it portable across any hardware that
supports population count and integer comparison.  The newer geometric
lane is deliberately narrower: it requires efficient floating-point
vector arithmetic only for the candidate set that survives the Boolean
filter.
